% Generated draft: requires mathematical and source review.
\documentclass[11pt]{article}
\usepackage[margin=25mm]{geometry}
\usepackage{fontspec}
\setmainfont{Latin Modern Roman}
\usepackage{amsmath,amssymb,mathtools,unicode-math}
\setmathfont{Latin Modern Math}
\usepackage{enumitem,needspace,xurl,hyperref}
\setlength{\parindent}{0pt}
\setlength{\parskip}{5pt}
\setlength{\emergencystretch}{3em}
\newcommand{\N}{\mathbb{N}}
\newcommand{\Z}{\mathbb{Z}}
\newcommand{\Q}{\mathbb{Q}}
\newcommand{\R}{\mathbb{R}}
\newcommand{\C}{\mathbb{C}}
\newcommand{\F}{\mathbb{F}}
\newcommand{\A}{\mathbb{A}}
\newcommand{\PP}{\mathbb{P}}
\newcommand{\E}{\mathbb{E}}
\newcommand{\eps}{\varepsilon}
\newcommand{\dd}{\,\mathrm{d}}
\begin{document}
\section*{UnsolvedMath ID: 800684}
\textbf{Source problem code:} catalogue entry
\subsection{Definitions and mathematical statement}
Let \(\mathbb{P}^{1}\) denote the complex projective line, and let \(\omega\) be its Fubini--Study form. Write \(\mathbb{D}(t)=\{z\in\mathbb{C}:|z|<t\}\). Whenever a holomorphic map \(h\) into \(\mathbb{P}^{1}\) is defined on a disc containing \(\mathbb{D}(r)\), its characteristic in the supplied convention is
\[
T_h(r):=\int_0^r \frac{\mathrm{d}t}{t}\int_{\mathbb{D}(t)}h^*(\omega),
\]
where \(h^*(\omega)\) is the pullback form and the integral at zero is interpreted as an improper integral.

\textbf{Problem (Dinh--Sibony, Problem 11, 2018).} Find the minimal growth of \(T_h\) for \(h\in\mathscr{U}\).

\textbf{Unresolved qualifications.} The supplied record does not define \(\mathscr{U}\), specify the domains of its maps, or state the limiting regime for \(r\). The conditional explanation of the integral above does not supply these missing hypotheses. The record also does not specify the comparison relation intended by ``minimal growth,'' or whether a minimum must be attained by a member of \(\mathscr{U}\). Consequently, the intended extremal problem cannot be stated with complete quantifiers from this record alone.
\subsection{Short English statement}
Dinh and Sibony ask for the least possible growth of the displayed Nevanlinna characteristic among maps in \(\mathscr{U}\). The class \(\mathscr{U}\), the radius limit, and the meaning of least growth are missing from the supplied record, so a precise extremal formulation requires further source context.
\subsection{Sources}
\begin{itemize}
\item Local UnsolvedMath registry, permanent ID \(800684\), exact source code \texttt{catalogue entry}, title ``Holomorphic Foliation Theory --- Problem 11.'' The registry attributes the problem to Dinh and Sibony, with proposed year 2018, and labels it open; this is a supplied status label, not an independently verified mathematical or literature assessment.
\item Dinh and Sibony, \emph{Some Open Problems on Holomorphic Foliation Theory} (2018), Problem 11. Supplied source URL: \url{https://arxiv.org/abs/1807.05640}. The record reports extraction by \texttt{arxiv\_latex} and an access date of 2026-10-16. These are supplied provenance metadata; the source has not been independently accessed or verified for this draft.
\end{itemize}
\end{document}
